% Concentration des ions cuivre(II) d'une pile Daniell qui débite un courant constant (exercice 8, chapitre 8)
\begin{tikzpicture}
\begin{axis}[otMesures,otAxesBout,otLecture,width=96mm,height=50mm,xmin=0,xmax=6.2,ymin=0,ymax=0.12,
  xtick={0,1,...,6},ytick={0,0.02,...,0.12},
  yticklabel style={/pgf/number format/fixed,/pgf/number format/fixed zerofill,/pgf/number format/precision=2},
  xlabel={$t$ (h)},ylabel={$[\ce{Cu^2+}]$ (mol$\cdot$L$^{-1}$)}]
\solution{\addplot[otGreenDark,line width=0.8pt,domain=0:5.35,samples=2] {0.1000-0.0187*x};
  \draw[otOrange,line width=0.6pt,dash pattern=on 2pt off 1.5pt] (axis cs:5.35,0) -- (axis cs:5.35,0.035);
  \node[pcLegende,anchor=south] at (axis cs:5.35,0.035) {5,3~h};
  \node[pcLegende,text=otGreenDark,anchor=south west] at (axis cs:2.3,0.066) {pente $a\approx\qty{-0,0187}{mol.L^{-1}.h^{-1}}$};}
\addplot[only marks,mark=*,mark size=1.8pt,otNavy] coordinates {(0,0.100)(1,0.081)(2,0.063)(3,0.044)(4,0.025)};
\end{axis}
\end{tikzpicture}
\par\smallskip\legende{Concentration des ions cuivre(II) mesurée toutes les heures pendant le fonctionnement de la pile.}
